\chapter{Graphs}

\begin{center}\itshape
A connected multigraph has a closed walk using every edge exactly once if and only if all its degrees are even.\\
\upshape(Euler, 1736; Hierholzer, 1873, in our words)
\end{center}

K\"onigsberg, in Prussia, was built on both banks of the river Pregel and on two islands, joined by seven bridges. Its
citizens, the story goes, liked to wonder on their Sunday walks whether one could cross every bridge exactly once and come
back home. In 1736 the question reached Leonhard Euler in St Petersburg. He thought at first that it had nothing to do with
mathematics, then saw why it did: the answer does not depend on distances or shapes, only on which land masses each bridge
joins, a ``geometry of position'' that Leibniz had wished for. Each land mass is touched by an odd number of bridges, and a
walk that enters and leaves a region uses its bridges in pairs; so no such walk exists. Euler's paper is the first theorem
of graph theory, although it contains no graph: the drawing of Figure~\ref{fig:konigsberg}, with points and lines, came
150 years later. The converse, that even degrees suffice, Euler stated without proof; Carl Hierholzer proved it, and died in
1871 before his paper could appear.

\begin{figure}[ht]
\centering
\fig{ch05-konigsberg}
\caption{The seven bridges of K\"onigsberg as a multigraph: four land masses (the two banks $N$ and $S$, the island $I$ and the
eastern part $E$), seven edges. The degrees are $3,3,5,3$, all odd, so there is no walk crossing each bridge once.}
\label{fig:konigsberg}
\end{figure}

Graphs are the most versatile objects of computer science. A road map, the web, a social network, the dependencies of a
build, the states of a protocol, the molecules of Chapter~3 and the conflicts between exams in a timetable are graphs. This
chapter lays the foundations used by the rest of the book: walks and connectivity, Euler's theorem and its most elegant
application, the de Bruijn sequences that read every word of length $n$ in a single circular word; the much harder
Hamiltonian cycles, with Dirac's theorem and R\'edei's theorem on tournaments; and Menger's theorem on disjoint paths, the
mother of the max-flow min-cut theorem of Chapter~6, together with its consequence for one-way streets.

\begin{think}
A safe has a keypad with digits $0$ to $9$ and opens as soon as the last four digits typed form the code, whatever came
before. A burglar who types all $10^4$ codes one after the other presses $40\,000$ keys. How few presses could suffice?
\end{think}

\section{Walks, paths and connectivity}

A \emph{graph} $G=(V,E)$ is a finite set $V$ of vertices with a set $E$ of edges, each an unordered pair $\{u,v\}$ of distinct vertices, written $uv$. In a
\emph{multigraph} several edges may join the same pair, and \emph{loops} (edges from a vertex to itself) are allowed. In a \emph{directed graph} the edges
(\emph{arcs}) are ordered pairs $(u,v)$, written $u\to v$. The \emph{degree} $\deg v$ is the number of edge ends at $v$ (a loop counts twice); in a directed
graph, $\deg^+v$ and $\deg^-v$ count the arcs leaving and entering $v$. A \emph{walk} is a sequence $v_0e_1v_1\cdots e_kv_k$ in which each edge $e_i$ joins $v_{i-1}$ to $v_i$
(in this direction, for a directed graph); it is \emph{closed} if $v_k=v_0$, a \emph{trail} if its edges are distinct, a \emph{path} if its vertices are distinct. A \emph{cycle} is
a closed trail whose vertices $v_1,\dots,v_k$ are distinct, $k\ge1$ for multigraphs (a loop, or two parallel edges, is a short cycle), $k\ge3$ for
graphs. The graph is \emph{connected} if any two vertices are joined by a walk; its maximal connected subgraphs are its \emph{components}.

\begin{lemma}\label{lem:walkpath}
A walk from $u$ to $v$ contains a path from $u$ to $v$ (using a subset of its edges). A closed walk of odd length contains an
odd cycle.
\end{lemma}
\begin{proof}
If a vertex repeats, $v_i=v_j$ with $i<j$, delete the part between them; the walk gets shorter, and after finitely many steps it is a
path. For the second claim, induct on the length $\ell$. If no vertex repeats except $v_0=v_\ell$, the walk is an odd cycle (for $\ell=1$, a
loop). Otherwise some $v_i=v_j$ with $0\le i<j\le\ell$ and $(i,j)\ne(0,\ell)$ splits the walk into two shorter closed walks, of lengths $j-i$ and
$\ell-(j-i)$, one of which is odd.
\end{proof}

A graph is \emph{bipartite} if its vertices can be split into two classes with every edge between the classes: jobs and workers,
students and schools, the two sides of a matching problem in Chapter~6.

\begin{theorem}[K\H onig, 1936]\label{thm:bipartite}
A graph is bipartite if and only if it has no cycle of odd length.
\end{theorem}
\begin{proof}
Along a cycle in a bipartite graph the classes alternate, so its length is even. Conversely, it suffices to treat a connected
graph. Fix a vertex $r$, let $d(v)$ be the length of a shortest walk from $r$ to $v$, and put $v$ in class $0$ or $1$ according to the parity of $d(v)$.
If an edge $uv$ joined two vertices of the same class, a shortest walk to $u$, the edge $uv$ and a shortest walk from $v$ back to $r$
would form a closed walk of odd length $d(u)+1+d(v)$, which contains an odd cycle by Lemma~\ref{lem:walkpath}.
\end{proof}

The proof is an algorithm: breadth-first search computes $d$ in time $O(|V|+|E|)$ and either returns the two classes or finds an
odd cycle, a certificate that none exist. This pattern, an algorithm that outputs either the object or a short proof of its
non-existence, will come back with every good theorem of this book.

\begin{proposition}[trees]\label{prop:treechar}
For a graph $T$ with $n$ vertices, the following are equivalent: (1) $T$ is a tree (connected, without cycles); (2) any two
vertices are joined by exactly one path; (3) $T$ is connected, and removing any edge disconnects it; (4) $T$ is connected and
has $n-1$ edges; (5) $T$ has no cycle and $n-1$ edges. Every connected graph contains a \emph{spanning tree}, a tree with all its vertices.
\end{proposition}
\begin{proof}
(1)$\Rightarrow$(2): two different paths from $u$ to $v$ would contain a cycle: follow the first until it leaves the second, then continue
along the first until it meets the second again, and return along the second. (2)$\Rightarrow$(3): if $T-uv$ were connected, a path from $u$ to $v$
in it and the edge $uv$ would be two paths. (3)$\Rightarrow$(1): an edge on a cycle can be removed without disconnecting. So (1), (2), (3) are
equivalent, and (1) implies (4) and (5) by Lemma~\ref{lem:trees}. A connected graph contains a spanning tree: remove edges lying
on cycles while there are any; connectivity is kept. If (4) holds, a spanning tree has $n-1$ edges, so it is $T$ itself. If (5) holds,
each component $C_i$ of $T$ is a tree with $n_i-1$ edges; summing, $n-1=n-c$ where $c$ is the number of components, so $c=1$.
\end{proof}

\section{Euler tours}

An \emph{Euler tour} of a multigraph is a closed trail using every edge.

\begin{theorem}[Euler, Hierholzer]\label{thm:euler}
A connected multigraph has an Euler tour if and only if every vertex has even degree. A directed multigraph in which
every vertex with an arc is reachable from every other such vertex has an Euler tour if and only if $\deg^+v=\deg^-v$ for every $v$.
\end{theorem}
\begin{proof}
Necessity: each passage of the tour through a vertex uses two edge ends, one in and one out, and so does the start and
end of the tour; a loop contributes $2$ to one passage.

Sufficiency, undirected. First, if all degrees are even, every edge lies on a closed trail: starting from it, walk along unused
edges; when the walk enters a vertex $w\ne v_0$, it has used an odd number of the edge ends at $w$, so an unused one remains, and the
walk can only stop at $v_0$. Now let $W$ be a longest closed trail. If it misses an edge, by connectivity some missed edge $e$ is at a vertex
$w$ of $W$ (take a path from a missed edge to $W$; its last edge before reaching $W$ is missed, or the missed edge is at $W$ already).
Removing the edges of $W$ leaves all degrees even, so $e$ lies on a closed trail $W'$ disjoint in edges from $W$. Splicing $W'$ into $W$ at
$w$ gives a longer closed trail, a contradiction. The directed case is the same, with ``enters by an arc, leaves by an unused arc''.
\end{proof}

The proof is Hierholzer's algorithm, which runs in linear time with a stack: follow unused edges until stuck, then back up,
emitting the edges in reverse; every backtracking point starts a detour that gets spliced in. Corollary: a connected multigraph
has an Euler \emph{trail} (not closed) if and only if it has exactly two vertices of odd degree, its ends: add an edge between them.
Routing a snowplough, a street sweeper or a printer head through every segment of a drawing once is this theorem
(Exercise~\ref{ex:trails} handles more odd vertices).

\section{De Bruijn sequences}

Now the safe. A cyclic word $c_0c_1\cdots c_{L-1}$ over an alphabet of $q$ letters is a \emph{de Bruijn sequence of order $n$} if each of the $q^n$ words of length
$n$ appears exactly once among its $L$ windows $c_ic_{i+1}\cdots c_{i+n-1}$ (indices mod $L$); then $L=q^n$. For $q=2$, $n=3$: $00010111$, whose windows are $000$, $001$,
$010$, $101$, $011$, $111$, $110$, $100$. Unrolled, it gives a linear word of length $q^n+n-1$ containing every word of length $n$, the shortest possible,
since each of the $q^n$ words must end at a different position.

The \emph{de Bruijn graph} $D(q,n)$ has the $q^{n-1}$ words of length $n-1$ as vertices, and an arc $a_1\cdots a_{n-1}\to a_2\cdots a_n$ labelled $a_1\cdots a_n$ for every word of length
$n$: one shifts in a letter on the right and drops one on the left (Figure~\ref{fig:debruijn}). Loops are allowed: $00\to00$ is the word $000$.

\begin{figure}[ht]
\centering
\fig{ch05-debruijn}
\caption{The de Bruijn graph $D(2,4)$: the vertices are the $8$ binary words of length $3$, and the $16$ arcs the words of length $4$; each
vertex has two arcs in and two out. An Euler tour, read by the first letter of each arc, is a de Bruijn sequence of order $4$,
such as $0000100110101111$.}
\label{fig:debruijn}
\end{figure}

\begin{theorem}[de Bruijn, 1946; Flye Sainte-Marie, 1894]\label{thm:debruijn}
De Bruijn sequences exist for all $q\ge1$ and $n\ge1$.
\end{theorem}
\begin{proof}
For $n=1$ write the $q$ letters. For $n\ge2$, in $D(q,n)$ every vertex has $q$ arcs out (choose the letter shifted in) and $q$ in (choose the
letter dropped), and every vertex $b_1\cdots b_{n-1}$ is reachable from every vertex $a_1\cdots a_{n-1}$, by shifting in $b_1,\dots,b_{n-1}$. By
Theorem~\ref{thm:euler} there is an Euler tour $e_0e_1\cdots e_{L-1}$, $L=q^n$. Let $c_i$ be the first letter of $e_i$. Consecutive arcs overlap: if $e_i$ is $a_1\cdots a_n$, then
$e_{i+1}=a_2\cdots a_nb$, so the first letter of $e_{i+j}$ is $a_{j+1}$ for $0\le j<n$, by induction on $j$. Hence the window $c_ic_{i+1}\cdots c_{i+n-1}$ is $e_i$, and since the tour uses
every arc once, every word of length $n$ is a window exactly once.
\end{proof}

So the burglar needs $10^4+3=10\,003$ key presses, not $40\,000$. De Bruijn sequences are used to locate a position from a short
window (a rotary encoder, a robot reading a coded strip on the floor, a camera reading a structured-light pattern), to test every
input of a shift register, and in the card trick of Exercise~\ref{ex:cardtrick}. How many are there? Camille Flye Sainte-Marie
counted them in 1894, in a problem in the \emph{Interm\'ediaire des math\'ematiciens} that nobody noticed: for $q=2$ there are $2^{2^{n-1}-n}$ cyclic
sequences. De Bruijn rediscovered the count in 1946, and the general formula $(q!)^{q^{n-1}}/q^n$ follows from the BEST theorem, which counts the
Euler tours of a directed graph (Chapter~9).

\begin{remark}
The lexicographically smallest de Bruijn sequence has a beautiful description (Fredricksen and Maiorana, 1978): concatenate,
in lexicographic order, the Lyndon words whose length divides $n$ (Exercise~\ref{ex:lyndon}). For $q=2$, $n=4$: $0\,0001\,0011\,01\,0111\,1$. The
notebook checks it, and generates the sequences in constant amortised time per letter.
\end{remark}

\section{Hamiltonian cycles and tournaments}

A \emph{Hamiltonian cycle} visits every \emph{vertex} exactly once. The question looks like Euler's, with vertices instead of edges, yet it is of a
completely different nature: no simple criterion is known, and Tome~C3 will show that deciding it is NP-complete, so that a
fast algorithm would solve thousands of other problems at once. The name comes from William Rowan Hamilton's \emph{Icosian game} of
1857, a puzzle sold in London that asked for such a cycle on the dodecahedron; it did not sell well. What one can prove are
sufficient conditions, the first and most famous being Dirac's.

\begin{theorem}[Dirac, 1952]\label{thm:dirac}
A graph with $n\ge3$ vertices in which every vertex has degree at least $\frac n2$ has a Hamiltonian cycle.
\end{theorem}
\begin{proof}
The graph is connected: two non-adjacent vertices have together at least $n$ neighbours among the $n-2$ others, so a common one. Let
$P=v_0v_1\cdots v_k$ be a longest path. All neighbours of $v_0$ and of $v_k$ lie on $P$, otherwise $P$ could be extended. Let $I=\{i:v_0v_{i+1}\in E\}$ and
$J=\{i:v_iv_k\in E\}$, subsets of the $k$-element set $\{0,\dots,k-1\}$ of sizes $\ge\frac n2$. As $|I|+|J|\ge n>k$, they meet. Take $i\in I\cap J$: then
\[
v_0v_1\cdots v_i\,v_kv_{k-1}\cdots v_{i+1}\,v_0
\]
is a cycle $C$ through the $k+1$ vertices of $P$. If $k+1<n$, a vertex $u$ outside $C$ has a neighbour on $C$ (connectivity); starting at $u$
and going around $C$ gives a path with $k+2$ vertices, contradicting the maximality of $P$. So $C$ is Hamiltonian.
\end{proof}

The bound is sharp: two cliques of $\frac{n+1}2$ vertices sharing one vertex have minimum degree $\frac{n-1}2$ and no Hamiltonian cycle
(Exercise~\ref{ex:diracsharp}). The proof is again an algorithm, the \emph{rotation} trick, which is the basis of the best heuristics
for the travelling salesman problem.

\subsection*{Tournaments}

In a round-robin tournament every two players meet once and there are no draws: a \emph{tournament} is a directed graph with exactly
one arc between any two vertices. Can the players always be ranked so that each one beat the next? The answer is yes,
whatever the results.

\begin{theorem}[R\'edei, 1934]\label{thm:redei}
Every tournament has a Hamiltonian path.
\end{theorem}
\begin{proof}
By induction on the number of players, the case of one player being trivial. Remove a player $x$ and rank the others as a path
$v_1\to v_2\to\cdots\to v_m$. If $x\to v_1$, put $x$ first; if $v_m\to x$, put it last. Otherwise $v_1\to x$ and $x\to v_m$, so there is a first index $i$ with
$x\to v_{i+1}$, and then $v_i\to x$ (as $i$ is the first, and $v_1\to x$): insert $x$ between $v_i$ and $v_{i+1}$.
\end{proof}

This is insertion sort, run on a ``comparison'' that need not be transitive: with binary search for the insertion point
(the first and last players play the role of the ends of the interval) it finds a Hamiltonian path with $O(n\log n)$ games looked
up. R\'edei proved more: the number of Hamiltonian paths is always \emph{odd}, so a ranking always exists. It is unique only when the
tournament is transitive: when everybody beat everybody below them in the
pecking order, the order is the ranking; when there are cycles, being ranked first means little (Exercise~\ref{ex:camion} shows
that strongly connected tournaments have Hamiltonian \emph{cycles}, so that every player can be put first). Landau, studying the
pecking orders of hens in 1951, proved that some player is a \emph{king}: it beats every other player directly or through an
intermediary (Exercise~\ref{ex:king}).

\section{Connectivity and Menger's theorem}

How robust is a network? A graph is \emph{$k$-connected} if it has more than $k$ vertices and stays connected after the removal of any $k-1$
vertices. For two sets of vertices $A$ and $B$, an \emph{$A$--$B$ path} is a path whose first vertex is in $A$, whose last vertex is in $B$, and whose
other vertices are in neither (a single vertex of $A\cap B$ is such a path). A set $X$ of vertices \emph{separates} $A$ from $B$ if every $A$--$B$ path
contains a vertex of $X$; for instance $A$ and $B$ themselves do. Paths are \emph{disjoint} if they share no vertex. Clearly the number of
disjoint $A$--$B$ paths is at most the size of any separator, since each path uses a different vertex of it. Karl Menger proved in
1927 that this obvious obstruction is the only one. His proof had a gap, filled by D\'enes K\H onig in 1931; the following short
proof is due to Frank G\"oring (2000).

\begin{theorem}[Menger, 1927]\label{thm:menger}
Let $G$ be a graph and $A,B\subseteq V$. The maximum number of disjoint $A$--$B$ paths equals the minimum size of a set separating $A$ from $B$.
\end{theorem}
\begin{proof}
Let $k$ be the minimum size of a separator; we show that there are $k$ disjoint $A$--$B$ paths, by induction on the number of edges.
If $G$ has no edge, the $A$--$B$ paths are the vertices of $A\cap B$, and $A\cap B$ is the smallest separator. Otherwise take an edge $e=xy$, and
let $G/e$ be the graph in which $x$ and $y$ are merged into one new vertex $v_e$ (adjacent to all neighbours of $x$ and $y$), with $A$ and $B$
modified accordingly ($v_e$ belongs to $A$ if $x$ or $y$ does, and similarly for $B$).

If $G/e$ has $k$ disjoint $A$--$B$ paths, so does $G$: replace $v_e$, if a path uses it, by $x$, $y$ or the edge $xy$, as needed. Otherwise, by the
induction hypothesis (as $G/e$ has fewer edges), $G/e$ has a separator $Y$ with fewer than $k$ vertices. Then $v_e\in Y$, since otherwise $Y$ would
separate $A$ from $B$ in $G$. So $X=(Y\setminus\{v_e\})\cup\{x,y\}$ separates $A$ from $B$ in $G$, and $|X|\le k$, hence $|X|=k$.

Now consider $G-e$. Every set $S$ separating $A$ from $X$ in $G-e$ separates $A$ from $B$ in $G$: an $A$--$B$ path $P$ in $G$ meets $X$, and its initial
segment up to the first vertex of $X$ is an $A$--$X$ path that does not use $e$ (whose ends are both in $X$), so it meets $S$. Hence such an $S$ has at
least $k$ vertices, and by induction ($G-e$ has fewer edges) there are $k$ disjoint $A$--$X$ paths in $G-e$; as $|X|=k$, each ends at a different
vertex of $X$. Likewise there are $k$ disjoint $X$--$B$ paths, one starting at each vertex of $X$. An $A$--$X$ path and an $X$--$B$ path meet only in $X$:
if they shared a vertex outside $X$, their union would contain an $A$--$B$ path avoiding $X$. So joining them at the vertices of $X$ gives $k$
disjoint $A$--$B$ walks, which are paths.
\end{proof}

\begin{corollary}\label{cor:mengerst}
If $s$ and $t$ are non-adjacent vertices, the maximum number of $s$--$t$ paths with no common vertex other than $s$ and $t$ equals the
minimum number of vertices (other than $s$, $t$) whose removal separates $s$ from $t$. Consequently a graph with more than $k$
vertices is $k$-connected if and only if any two non-adjacent vertices are joined by $k$ such paths.
\end{corollary}
\begin{proof}
Apply the theorem in $G-s-t$ with $A$ the neighbours of $s$ and $B$ those of $t$: disjoint $A$--$B$ paths extend to $s$--$t$ paths meeting only at
their ends, and conversely, and a separator of $A$ from $B$ in $G-s-t$ is exactly a set of vertices whose removal separates $s$ from $t$. For
the consequence: if the graph is $k$-connected, separating two non-adjacent vertices takes $k$ vertices, so there are $k$ paths;
conversely, if removing a set $S$ of fewer than $k$ vertices disconnected the graph, two vertices in different components would
be non-adjacent and separated by $S$. (If all pairs are adjacent, the graph is complete, and $k$-connected as it has more than $k$
vertices.)
\end{proof}

Chapter~6 proves the version for edges and for capacities, the max-flow min-cut theorem, with an algorithm; and Hall's
marriage theorem and K\H onig's theorem on bipartite matchings are Menger's theorem in disguise
(Exercise~\ref{ex:fan} is another useful form).

\subsection*{Ears and one-way streets}

The $2$-connected graphs have a beautiful inductive structure. An \emph{ear} of a subgraph $H$ of $G$ is a path of $G$ whose two ends are distinct
vertices of $H$ and whose other vertices and edges are not in $H$; a single edge of $G$ between two vertices of $H$, not in $H$, is an ear. An
\emph{ear decomposition} of $G$ is a sequence $H_0\subseteq H_1\subseteq\cdots\subseteq H_m=G$ where $H_0$ is a cycle and $H_{i+1}$ is $H_i$ plus an ear of $H_i$ (Figure~\ref{fig:ears}).

\begin{figure}[ht]
\centering
\fig{ch05-ears}
\caption{An ear decomposition: a cycle (thick), then three ears added one after the other (numbered), each a path between two
distinct vertices already reached. Orienting the cycle and each ear in one direction gives a strongly connected orientation
(Theorem~\ref{thm:robbins}).}
\label{fig:ears}
\end{figure}

\begin{theorem}[Whitney, 1932]\label{thm:ears}
A graph with at least three vertices is $2$-connected if and only if it has an ear decomposition.
\end{theorem}
\begin{proof}
A cycle is $2$-connected, and adding an ear $P$ between distinct vertices $u,v$ of a $2$-connected $H$ keeps it $2$-connected: removing a vertex of
$H$ leaves $H$ connected and every vertex of $P$ still attached to $u$ or $v$ along $P$; removing an inner vertex of $P$ leaves $H$ intact and
the two pieces of $P$ attached to $u$ and to $v$.

Conversely, let $G$ be $2$-connected. Every vertex has degree at least $2$ (one neighbour could be removed to isolate it), so $G$
contains a cycle (the proof of Lemma~\ref{lem:trees}(1)). Let $H$ be a maximal subgraph of $G$ with an ear decomposition starting from a
cycle. If an edge of $G$ joins two vertices of $H$ and is not in $H$, it is an ear, contradicting maximality. So if $H\ne G$, some vertex is
outside $H$ and, by connectivity, an edge $uv$ has $u\in H$, $v\notin H$. As $G-u$ is connected and $H$ has a vertex other than $u$, there is a path
in $G-u$ from $v$ to $H$; stop it at its first vertex $w$ of $H$. Then $u$, $v$, the path to $w$ form an ear of $H$, again a contradiction.
So $H=G$.
\end{proof}

In 1939 Herbert Robbins, then a topologist, answered a question of town planners: when can all the streets of a town be made
one-way so that one can still drive from anywhere to anywhere? A \emph{bridge} is an edge whose removal disconnects the graph; a bridge
obviously cannot be made one-way. That is the only obstruction.

\begin{theorem}[Robbins, 1939]\label{thm:robbins}
A connected graph has a strongly connected orientation if and only if it has no bridge.
\end{theorem}
\begin{proof}
If $uv$ is a bridge oriented $u\to v$, nothing leads back from $v$ to $u$. Conversely, let $G$ be connected without bridges, with at least two
vertices. Let $H$ be a maximal subgraph with a strongly connected orientation in which every edge joining two vertices of $H$ belongs to
$H$; a cycle exists (an edge on no cycle is a bridge) and, with its chords, oriented along the cycle for the cycle edges and
arbitrarily for the chords, is such a subgraph. If $H\ne G$, take an edge $uv$ with $u\in H$, $v\notin H$. As $uv$ is not a bridge, $G-uv$ has a path from
$v$ to $u$; stop it at its first vertex $w$ in $H$ (possibly $w=u$). Orient $u\to v\to\cdots\to w$: every new vertex can be reached from $u$ and reaches $w$,
hence everything; add the new vertices and all remaining edges between vertices of the enlarged $H$, oriented arbitrarily. This
contradicts maximality, so $H=G$.
\end{proof}

A depth-first search finds such an orientation in linear time: orient the tree edges away from the root and the other edges
towards the root (Exercise~\ref{ex:robbinsdfs}).

\begin{beauty}
Euler's theorem, Menger's theorem and Robbins's theorem share a pattern that runs through this book: an obstruction that is
obviously fatal (an odd degree, a small separator, a bridge) turns out to be the only one. The proofs are constructive, and the
obstruction is a certificate: whoever claims that no tour, no $k$ disjoint paths, no good orientation exists can prove it in a line.
\end{beauty}

\section{Exercises}

\begin{exercise}[$\star$]\label{ex:trails}
Show that the edges of a connected graph with $2k>0$ vertices of odd degree can be partitioned into $k$ trails, and not fewer.
\end{exercise}

\begin{exercise}[$\star$]\label{ex:dominoes}
A set of dominoes contains one tile for each unordered pair $\{a,b\}$ with $0\le a\le b\le6$, doubles included: $28$ tiles. Show that all of them can
be laid in a closed ring, with touching ends equal. What about the set with $0\le a\le b\le7$?
\end{exercise}

\begin{exercise}[$\star$]\label{ex:cardtrick}
\emph{A mind-reading trick.} Arrange $32$ cards in a circle so that any five consecutive cards are determined by their five colours
(red or black). Explain how a magician, told only the colours of the five cards that five spectators cut from the deck,
names the cards. How many arrangements of the colours are possible, up to rotation, for $16$ cards and windows of $4$?
\end{exercise}

\begin{exercise}[$\star\star$]\label{ex:camion}
\emph{Camion, 1959.} A tournament is \emph{strongly connected} if every player reaches every other along arcs. Show that a strongly
connected tournament with $n\ge3$ players has a Hamiltonian cycle.
\end{exercise}

\begin{exercise}[$\star$]\label{ex:king}
\emph{Landau's kings.} Show that in every tournament, a player with the most wins reaches every other player by a path of
length at most $2$.
\end{exercise}

\begin{exercise}[$\star$]\label{ex:diracsharp}
Show that Dirac's bound cannot be lowered: for every $n\ge3$ there is a graph with minimum degree $\lceil\frac n2\rceil-1$ and no Hamiltonian cycle.
\end{exercise}

\begin{exercise}[$\star$]\label{ex:gray}
\emph{Gray codes.} Show that the $n$-dimensional cube, whose vertices are the words of $\{0,1\}^n$ joined when they differ in one bit, has a
Hamiltonian cycle for $n\ge2$, and give it explicitly. Why do rotary encoders and Karnaugh maps use it?
\end{exercise}

\begin{exercise}[$\star\star$]\label{ex:fan}
\emph{The fan lemma.} Let $G$ be $k$-connected, $v$ a vertex and $U$ a set of at least $k$ vertices not containing $v$. Show that there are $k$ paths from
$v$ to $U$, sharing only $v$, each with only its last vertex in $U$.
\end{exercise}

\begin{exercise}[$\star\star$]\label{ex:robbinsdfs}
Run a depth-first search on a connected bridgeless graph, orient the tree edges from parent to child and the other edges from
descendant to ancestor. Prove that the orientation is strongly connected.
\end{exercise}

\begin{exercise}[$\star\star$]\label{ex:linegraph}
The \emph{line graph} $L(G)$ has the edges of $G$ as vertices, two of them adjacent when they share an end. Show that if $G$ has an Euler
tour then $L(G)$ has a Hamiltonian cycle. Is the converse true?
\end{exercise}

\begin{exercise}[$\star\star$]\label{ex:twoconnected}
Show that in a $2$-connected graph any two vertices lie on a common cycle, and any two edges too.
\end{exercise}

\section{Hints}

\begin{hint}{ex:trails}
Add $k$ new edges pairing the odd vertices.
\end{hint}
\begin{hint}{ex:dominoes}
The complete graph on $\{0,\dots,6\}$ with a loop at each vertex.
\end{hint}
\begin{hint}{ex:cardtrick}
Red is $1$, black $0$: a de Bruijn sequence of order $5$. For the count, Flye Sainte-Marie's formula.
\end{hint}
\begin{hint}{ex:camion}
Start from a cycle and absorb the other players one by one, or grow the longest cycle.
\end{hint}
\begin{hint}{ex:king}
If $y$ beats the champion $x$, look at the players beaten by $x$.
\end{hint}
\begin{hint}{ex:diracsharp}
Two cliques sharing a vertex, or a complete bipartite graph with unequal sides.
\end{hint}
\begin{hint}{ex:gray}
Reflect: list the cycle for $n-1$ with a $0$ in front, then the same list backwards with a $1$ in front.
\end{hint}
\begin{hint}{ex:fan}
Add a new vertex joined to all of $U$, or apply Menger with $A=N(v)$ and $B=U$.
\end{hint}
\begin{hint}{ex:robbinsdfs}
Every vertex reaches the root: for a subtree, some back edge leaves it, otherwise the edge into it is a bridge.
\end{hint}
\begin{hint}{ex:linegraph}
Consecutive edges of the tour share a vertex. For the converse, try a star.
\end{hint}
\begin{hint}{ex:twoconnected}
Two vertices: Menger. Two edges: subdivide them.
\end{hint}

\section{Solutions}

\begin{solution}{ex:trails}
\idea Each trail has two ends, and odd vertices must be ends; pairing them by extra edges reduces to Euler.
\writeup Each odd vertex is an end of some trail of any partition (a trail passing through a vertex uses its edges in pairs), and a trail has
two ends, so at least $k$ trails are needed. Add $k$ new edges pairing the $2k$ odd vertices: all degrees become even and the multigraph is
connected, so it has an Euler tour (Theorem~\ref{thm:euler}). Deleting the $k$ new edges cuts the tour into $k$ trails covering all edges
once (as $k>0$, the tour contains at least one new edge, and each deletion of an edge of a closed trail cuts off one more piece).
\end{solution}

\begin{solution}{ex:dominoes}
\idea A tile is an edge between its two numbers; a ring is an Euler tour.
\writeup Take the vertices $0,\dots,6$, an edge for each tile $\{a,b\}$, $a\ne b$, and a loop for each double. Each vertex has $6$ edges to the others and
one loop, degree $6+2=8$, even; the graph is connected. An Euler tour lists the tiles with touching ends equal. With $0,\dots,7$, each vertex
has degree $7+2=9$, odd: no ring exists (and $8$ odd vertices: at least $4$ separate chains, by Exercise~\ref{ex:trails}).
\end{solution}

\begin{solution}{ex:cardtrick}
\idea The circle of colours must show every pattern of five colours once: a de Bruijn sequence.
\writeup Choose a binary de Bruijn sequence of order $5$, of length $32$, and assign the $32$ cards to its positions; a cut of five consecutive cards
is a window, and each colour pattern occurs at exactly one position, which the magician knows (in practice by a memorised rule, as
the sequences given by a shift register are easy to compute mentally). For $q=2$, $n=4$, Flye Sainte-Marie's formula gives
$2^{2^3-4}=16$ cyclic sequences of length $16$.
\end{solution}

\begin{solution}{ex:camion}
\idea Build up: a strongly connected tournament contains a cycle, and a longest cycle can absorb any remaining player.
\writeup A strongly connected tournament with $n\ge3$ players has a cycle of length $3$: take $x$, let $O$ be the players it beats and $I$ those beating it;
both are non-empty (strong connectivity), and some $y\in O$ beats some $z\in I$ (otherwise nothing leads from $O\cup\{x\}$ to $I$), giving $x\to y\to z\to x$.
Let $C=v_1\to\cdots\to v_m\to v_1$ be a longest cycle and suppose $m<n$. If a player $u\notin C$ has both an arc from and an arc to $C$, then going around $C$
there is an $i$ with $v_i\to u\to v_{i+1}$, and $u$ can be inserted: a longer cycle. So every player outside $C$ either beats all of $C$ (set $B$) or loses
to all of $C$ (set $L$). Both are non-empty: if $L$ were empty, no arc would enter $B$ from outside $B$, so $C$ could not reach $B$; if $B$ were empty,
no arc would leave $L$. A path from a player of $L$ to $C$ must leave $L$, and its first arc out of $L$ cannot go to $C$, which beats $L$;
so it goes to $B$, and some $\ell\in L$ beats some $b\in B$. Then $v_1\to\ell\to b\to v_3\to\cdots\to v_m\to v_1$ (using $v_1\to\ell$ and $b\to v_3$) has $m+1$
vertices, a contradiction. (For $m=3$ the cycle is $v_1\ell bv_3$; the vertex $v_2$ is dropped and two are added.)
\end{solution}

\begin{solution}{ex:king}
\idea A champion who lost to $y$ beats many players; if all of them also lost to $y$, then $y$ would have more wins.
\writeup Let $x$ have the most wins and let $y$ be any player not beaten by $x$, so $y\to x$. If $x$ beats some $z$ who beats $y$, then $x\to z\to y$. Otherwise $y$ beats
everyone beaten by $x$, and $x$ too, so $y$ has more wins than $x$, a contradiction.
\end{solution}

\begin{solution}{ex:diracsharp}
\idea A vertex whose removal disconnects the graph, or an imbalance in a bipartite graph, forbids a Hamiltonian cycle.
\writeup For $n$ odd, $n=2m+1$, take two cliques on $m+1$ vertices sharing one vertex $c$: degrees are $m$ or $2m$, so the minimum is $m=\lceil\frac n2\rceil-1$,
and removing $c$ disconnects the graph, while removing one vertex from a Hamiltonian cycle leaves a path. For $n$ even, $n=2m$, take
the complete bipartite graph with sides of sizes $m-1$ and $m+1$: the minimum degree is $m-1=\lceil\frac n2\rceil-1$, and a Hamiltonian cycle would
alternate between the sides, forcing them to have equal sizes.
\end{solution}

\begin{solution}{ex:gray}
\idea Two copies of the cube of dimension $n-1$, glued along matching vertices: traverse one, cross, traverse the other backwards.
\writeup Let $\Gamma_1=(0,1)$ and $\Gamma_n=0\Gamma_{n-1},1\overline{\Gamma_{n-1}}$, the list $\Gamma_{n-1}$ with $0$ in front followed by the reversed list with $1$ in front. By
induction consecutive words differ in one bit: inside each half by induction, and at the junction the words $0w$ and $1w$, $w$ the last
word of $\Gamma_{n-1}$. The last word of $\Gamma_n$ is $1$ followed by the first word of $\Gamma_{n-1}$, which is $0\cdots0$, so it differs in one bit from the first word $0\cdots0$
of $\Gamma_n$: a Hamiltonian cycle for $n\ge2$. Explicitly, the $k$-th word is $k\oplus\lfloor k/2\rfloor$ (bitwise exclusive or). An encoder reading a Gray code can
never misread a position by more than one step when it is caught between two positions, since only one bit changes.
\end{solution}

\begin{solution}{ex:fan}
\idea Turn ``one vertex to many'' into ``many to many'' and use Menger.
\writeup Add a new vertex $u$ joined to every vertex of $U$; the new graph is still $k$-connected (removing fewer than $k$ vertices leaves $G$ minus those
connected, or $u$ removed). Apply Theorem~\ref{thm:menger} with $A=N(v)$, $B=U$ in $G-v$: a set separating $A$ from $B$ in $G-v$ separates $v$ from $U$, hence from $u$ if
$v$ and $u$ are non-adjacent, which they are; so with Corollary~\ref{cor:mengerst} applied to $v$ and $u$ in the new graph, there are $k$ paths from $v$ to $u$
sharing only their ends. Dropping $u$ and cutting each path at its first vertex in $U$ gives the fan.
\end{solution}

\begin{solution}{ex:robbinsdfs}
\idea The root reaches everything along tree edges; it remains to see that everything reaches the root, which back edges provide.
\writeup In a depth-first search tree of an undirected graph every non-tree edge joins a vertex to one of its ancestors. The root reaches every vertex
along tree edges. Suppose some vertex does not reach the root, and take such a vertex $v$ as close to the root as possible; $v$ is not the root,
let $p$ be its parent, which reaches the root. No vertex of the subtree $T_v$ of $v$ reaches outside $T_v$, otherwise it would reach some vertex
outside, which is an ancestor of $v$ (back edges go to ancestors), closer to the root, hence reaching the root. So no back edge leaves $T_v$,
and the only edge between $T_v$ and the rest is the tree edge $pv$: it is a bridge, a contradiction.
\end{solution}

\begin{solution}{ex:linegraph}
\idea Read the Euler tour as a cyclic list of edges.
\writeup If $e_1e_2\cdots e_m$ is an Euler tour, consecutive edges share a vertex, and so do $e_m$ and $e_1$; every edge appears once, so $e_1,\dots,e_m$ is a
Hamiltonian cycle of $L(G)$ (for $m\ge3$). The converse fails: the star with three edges has a line graph which is a triangle, Hamiltonian,
but its centre has odd degree $3$.
\end{solution}

\begin{solution}{ex:twoconnected}
\idea Two internally disjoint paths form a cycle; for edges, put a new vertex in the middle of each.
\writeup For vertices $u\ne v$: if they are non-adjacent, Corollary~\ref{cor:mengerst} with $k=2$ gives two paths meeting only at $u$ and $v$, whose union is a
cycle. If $uv$ is an edge, $G-uv$ is still connected (otherwise, as $G$ has at least three vertices, $u$ or $v$ would separate it), so a path
from $u$ to $v$ in $G-uv$ plus $uv$ is a cycle. For two edges, subdivide each by a new vertex; the graph stays $2$-connected (it has an ear
decomposition: subdivide the corresponding ear), and a cycle through the two new vertices uses both edges.
\end{solution}

\section*{Sources and further reading}

\begin{itemize}
\item L.~Euler, \emph{Solutio problematis ad geometriam situs pertinentis}, Commentarii Academiae Scientiarum Petropolitanae 8
(1736, published 1741), 128--140; C.~Hierholzer, \emph{\"Uber die M\"oglichkeit, einen Linienzug ohne Wiederholung und ohne
Unterbrechung zu umfahren}, Mathematische Annalen 6 (1873), 30--32. Both are translated in N.~L.~Biggs, E.~K.~Lloyd and
R.~J.~Wilson, \emph{Graph Theory 1736--1936}, Oxford University Press, 1976, a delightful history.
\item N.~G.~de Bruijn, \emph{A combinatorial problem}, Proceedings of the Koninklijke Nederlandse Akademie van Wetenschappen 49
(1946), 758--764; H.~Fredricksen and J.~Maiorana, \emph{Necklaces of beads in $k$ colors and $k$-ary de Bruijn sequences}, Discrete
Mathematics 23 (1978), 207--210.
\item G.~A.~Dirac, \emph{Some theorems on abstract graphs}, Proceedings of the London Mathematical Society 2 (1952), 69--81.
\item L.~R\'edei, \emph{Ein kombinatorischer Satz}, Acta Litterarum ac Scientiarum Szeged 7 (1934), 39--43; P.~Camion,
\emph{Chemins et circuits hamiltoniens des graphes complets}, Comptes Rendus de l'Acad\'emie des Sciences 249 (1959), 2151--2152.
\item K.~Menger, \emph{Zur allgemeinen Kurventheorie}, Fundamenta Mathematicae 10 (1927), 96--115; F.~G\"oring, \emph{Short proof
of Menger's theorem}, Discrete Mathematics 219 (2000), 295--296.
\item H.~Whitney, \emph{Non-separable and planar graphs}, Transactions of the AMS 34 (1932), 339--362; H.~E.~Robbins, \emph{A
theorem on graphs, with an application to a problem of traffic control}, American Mathematical Monthly 46 (1939), 281--283.
\item R.~Diestel, \emph{Graph Theory}, 5th ed., Springer, 2017 (free online edition), Chapters 1 and 3: the standard modern text.
\item J.~A.~Bondy and U.~S.~R.~Murty, \emph{Graph Theory}, Springer, 2008.
\end{itemize}
